\babbaru{DAFTAR LAMPIRAN}{%
\petunjuk{Centang kolom Terlampir bila dokumen sudah dilampirkan dan isi nomor halaman. Setiap lampiran diawali halaman judul lampiran.}{}
}

\begingroup\singlespacing\footnotesize\noindent
\begin{tabularx}{\textwidth}{|c|>{\raggedright\arraybackslash}X|c|c|c|c|}
\hline
\rowcolor{shadeheader}\centering\textbf{No} & \centering\textbf{Lampiran} & \centering\textbf{Sempro} & \centering\textbf{Sidang} & \centering\textbf{Hal.} & \centering\arraybackslash\textbf{\checkmark} \\\hline
1 & Form Kendali Bimbingan & WAJIB & WAJIB & & [ ] \\\hline
2 & Berita Acara Seminar Proposal & Tidak diperlukan & WAJIB & & [ ] \\\hline
3 & Kode Sumber dan Repositori & Opsional & WAJIB & & [ ] \\\hline
4 & Peta Kode untuk Code Walkthrough & Opsional & WAJIB & & [ ] \\\hline
5 & Dataset dan Dokumentasi Data & Opsional & WAJIB & & [ ] \\\hline
6 & Hasil Uji Similarity & Tidak diperlukan & WAJIB & & [ ] \\\hline
7 & Log Sesi AI & Tidak diperlukan & WAJIB bila berlaku & & [ ] \\\hline
8 & Bukti Riwayat Versi Laporan dan Kode & Tidak diperlukan & WAJIB & & [ ] \\\hline
9 & Daftar Riwayat Hidup & Tidak diperlukan & Opsional & & [ ] \\\hline
\end{tabularx}
\endgroup

\babbaru{LAMPIRAN 1. FORM KENDALI BIMBINGAN}{%
\petunjuk{Minimal 4 pertemuan per semester, ditandatangani pembimbing atau Kartu Kendali dari TASYS.}{\statusdua{\WAJIB}{\WAJIB}}
}
\noindent\placeholder{[Tempelkan dokumen pada halaman ini]}\par

\babbaru{LAMPIRAN 2. BERITA ACARA SEMINAR PROPOSAL}{%
\petunjuk{Beserta catatan perbaikan dari penguji dan bukti perbaikannya.}{\statusdua{\TIDAKPERLU}{\WAJIB}}
}
\noindent\placeholder{[Tempelkan dokumen pada halaman ini]}\par

\babbaru{LAMPIRAN 3. KODE SUMBER DAN REPOSITORI}{%
\petunjuk{Tautan repositori git, versi atau commit final, dan potongan kode penting. Kode lengkap cukup di repositori.}{\statusdua{\OPSIONAL}{\WAJIB}}
}
\noindent\placeholder{[Tempelkan dokumen pada halaman ini]}\par

\babbaru{LAMPIRAN 4. PETA KODE UNTUK CODE WALKTHROUGH}{%
\petunjuk{Wajib bila TA memuat implementasi perangkat lunak. Petakan modul utama agar penguji dapat menelusuri kode saat sidang.}{\statusdua{\OPSIONAL}{\WAJIB}}
\petunjuk{Uraikan fungsi, algoritma atau keputusan teknis, dan alasan pemilihan untuk setiap modul / berkas.}{}
}
\placeholder{[Tulis uraian dalam bentuk paragraf.]}\par

\babbaru{LAMPIRAN 5. DATASET DAN DOKUMENTASI DATA}{%
\petunjuk{Tautan dataset terbuka (Zenodo, Figshare, GitHub, repositori ITERA) beserta lisensi, atau metadata bila data terbatas (Pedoman TA 5.9).}{\statusdua{\OPSIONAL}{\WAJIB}}
}
\noindent\placeholder{[Tempelkan dokumen pada halaman ini]}\par

\babbaru{LAMPIRAN 6. HASIL UJI SIMILARITY}{%
\petunjuk{Laporan Turnitin lengkap (highlighted matches), similarity $\le$25\%, kutipan langsung $\le$10\%.}{\statusdua{\TIDAKPERLU}{\WAJIB}}
}
\noindent\placeholder{[Tempelkan dokumen pada halaman ini]}\par

\babbaru{LAMPIRAN 7. LOG SESI AI}{%
\petunjuk{Ekspor sesi AI yang relevan bila memakai AI. Boleh berupa tautan ke repositori digital.}{\statusdua{\TIDAKPERLU}{\WAJIBBILA}}
}
\noindent\placeholder{[Tempelkan dokumen pada halaman ini]}\par

\babbaru{LAMPIRAN 8. BUKTI RIWAYAT VERSI LAPORAN DAN KODE}{%
\petunjuk{Tangkapan layar riwayat commit dan riwayat versi dokumen, atau tautan aksesnya.}{\statusdua{\TIDAKPERLU}{\WAJIB}}
}
\noindent\placeholder{[Tempelkan dokumen pada halaman ini]}\par

\babbaru{LAMPIRAN 9. DAFTAR RIWAYAT HIDUP}{%
\petunjuk{CV singkat penulis.}{\statusdua{\TIDAKPERLU}{\OPSIONAL}}
}
\noindent\placeholder{[Tempelkan dokumen pada halaman ini]}\par
\clearpage
